% ============================================================================
% CHƯƠNG 3: KIẾN TRÚC FASTVIT
% ============================================================================
\chapter{Kiến trúc FastViT}
\label{chap:architecture}

Chương này phân tích chi tiết kiến trúc \fastvit{} \cite{fastvit2023}, bao gồm tổng quan thiết kế, các thành phần cốt lõi, cơ chế Structural Reparameterization, và so sánh các phiên bản.

% ============================================================================
\section{Tổng quan FastViT}
\label{sec:fastvit_overview}

\fastvit{} là một mô hình lai (hybrid) kết hợp CNN và Transformer, được thiết kế theo triết lý:

\begin{quote}
\textit{``Sử dụng token mixer hiệu quả nhất cho từng mức độ phân giải: RepMixer (dựa trên depthwise convolution) cho các stage có độ phân giải cao, Self-Attention cho stage có độ phân giải thấp.''}
\end{quote}

Điểm đột phá của \fastvit{} nằm ở việc áp dụng \textbf{Structural Reparameterization} xuyên suốt kiến trúc, cho phép mô hình sử dụng kiến trúc multi-branch phức tạp khi huấn luyện (tăng khả năng biểu diễn) nhưng rút gọn về chuỗi các phép tích chập đơn giản khi suy luận (tối ưu tốc độ).

% ============================================================================
\section{Kiến trúc tổng thể}
\label{sec:overall_architecture}

Kiến trúc tổng thể của \fastvit{} được tổ chức theo mô hình phân cấp (hierarchical) 4 stage, tương tự ResNet và Swin Transformer (Hình~\ref{fig:fastvit_arch}).

\begin{figure}[H]
\centering
\begin{tikzpicture}[
    block/.style={rectangle, draw, fill=blue!15, text width=4cm, 
                  text centered, rounded corners, minimum height=1cm, font=\small},
    repblock/.style={rectangle, draw, fill=green!15, text width=4cm,
                     text centered, rounded corners, minimum height=1cm, font=\small},
    attblock/.style={rectangle, draw, fill=orange!15, text width=4cm,
                     text centered, rounded corners, minimum height=1cm, font=\small},
    downblock/.style={rectangle, draw, fill=red!10, text width=4cm,
                      text centered, rounded corners, minimum height=0.8cm, font=\small},
    arrow/.style={->, thick},
    note/.style={font=\scriptsize, text=gray}
]
    % Input
    \node[block] (input) {Input Image\\$(B, 3, 256, 256)$};
    
    % Stem
    \node[block, below=0.4cm of input] (stem) {Convolutional Stem\\MobileOne Blocks, stride 4\\$(B, C_0, 64, 64)$};
    
    % Stage 0
    \node[repblock, below=0.4cm of stem] (stage0) {Stage 0: RepMixer Blocks\\$L_0$ blocks, dim $C_0$\\$(B, C_0, 64, 64)$};
    
    % Downsample 1
    \node[downblock, below=0.3cm of stage0] (down1) {Downsample: PatchEmbed\\stride 2};
    
    % Stage 1
    \node[repblock, below=0.3cm of down1] (stage1) {Stage 1: RepMixer Blocks\\$L_1$ blocks, dim $C_1$\\$(B, C_1, 32, 32)$};
    
    % Downsample 2
    \node[downblock, below=0.3cm of stage1] (down2) {Downsample: PatchEmbed\\stride 2};
    
    % Stage 2
    \node[repblock, below=0.3cm of down2] (stage2) {Stage 2: RepMixer Blocks\\$L_2$ blocks, dim $C_2$\\$(B, C_2, 16, 16)$};
    
    % Downsample 3
    \node[downblock, below=0.3cm of stage2] (down3) {Downsample: PatchEmbed\\stride 2};
    
    % Stage 3
    \node[attblock, below=0.3cm of down3] (stage3) {Stage 3: Attention Blocks\\$L_3$ blocks, dim $C_3$\\$(B, C_3, 8, 8)$};
    
    % Head
    \node[block, below=0.4cm of stage3] (head) {Classification Head\\GAP + MobileOne + Linear};
    
    % Arrows
    \draw[arrow] (input) -- (stem);
    \draw[arrow] (stem) -- (stage0);
    \draw[arrow] (stage0) -- (down1);
    \draw[arrow] (down1) -- (stage1);
    \draw[arrow] (stage1) -- (down2);
    \draw[arrow] (down2) -- (stage2);
    \draw[arrow] (stage2) -- (down3);
    \draw[arrow] (down3) -- (stage3);
    \draw[arrow] (stage3) -- (head);
    
    % Annotations
    \node[note, right=0.3cm of stage0] {Depthwise Conv};
    \node[note, right=0.3cm of stage1] {Depthwise Conv};
    \node[note, right=0.3cm of stage2] {Depthwise Conv};
    \node[note, right=0.3cm of stage3] {MHSA + RepCPE};
\end{tikzpicture}
\caption{Kiến trúc tổng thể của \fastvit{} (phiên bản SA). Các stage 0--2 sử dụng RepMixer (xanh lá), stage 3 sử dụng Attention (cam). Downsampling (đỏ nhạt) giữa các stage.}
\label{fig:fastvit_arch}
\end{figure}

% ----------------------------------------------------------------------------
\subsection{Input Stem}
\label{subsec:stem}

Input Stem là phần đầu tiên xử lý ảnh đầu vào, sử dụng chuỗi 3 MobileOne Blocks:

\begin{enumerate}
    \item \textbf{Conv $3 \times 3$, stride 2, groups=1:} Giảm kích thước không gian $\times 2$ và ánh xạ 3 kênh RGB $\to C_0$ kênh.
    \item \textbf{Conv $3 \times 3$, stride 2, groups=$C_0$:} Giảm kích thước $\times 2$ lần nữa (depthwise).
    \item \textbf{Conv $1 \times 1$, stride 1, groups=1:} Pointwise conv để tinh chỉnh biểu diễn kênh.
\end{enumerate}

Tổng cộng stem giảm kích thước $\times 4$: ảnh $256 \times 256 \to$ feature map $64 \times 64$ với $C_0$ kênh. Chi phí tính toán của stem (với $C_0 = 64$):

\begin{equation}
    \text{FLOPs}_{\text{stem}} \approx 3 \times 9 \times 64 \times 128^2 + 9 \times 64 \times 64^2 + 1 \times 64 \times 64 \times 64^2 \approx 46\text{M FLOPs}
\end{equation}

\noindent Đây là chi phí rất thấp, chiếm chưa đến 5\% tổng FLOPs của mô hình.

% ----------------------------------------------------------------------------
\subsection{Hierarchical Stages}
\label{subsec:stages}

\fastvit{} sử dụng 4 stages phân cấp với kích thước feature map giảm dần và số kênh tăng dần:

\begin{table}[H]
\centering
\caption{Kích thước feature map và token mixer tại mỗi stage (SA12, ảnh $256 \times 256$).}
\label{tab:stage_detail}
\begin{tabular}{c c c c c c}
\toprule
\textbf{Stage} & \textbf{Resolution} & \textbf{Channels} & \textbf{Blocks} & \textbf{Token Mixer} & \textbf{Tokens} \\
\midrule
0 & $64 \times 64$ & 64 & 2 & RepMixer & 4{,}096 \\
1 & $32 \times 32$ & 128 & 2 & RepMixer & 1{,}024 \\
2 & $16 \times 16$ & 256 & 6 & RepMixer & 256 \\
3 & $8 \times 8$ & 512 & 2 & Attention & 64 \\
\bottomrule
\end{tabular}
\end{table}

\textbf{Phân tích chiến lược thiết kế:}
\begin{itemize}
    \item \textbf{Stage 0--2 (RepMixer):} Feature map lớn $\Rightarrow$ nhiều token $\Rightarrow$ Self-Attention quá đắt. RepMixer (depthwise conv) có chi phí $\bigO{N \cdot k^2}$ tuyến tính theo số token, phù hợp.
    
    \item \textbf{Stage 3 (Attention):} Feature map nhỏ $8 \times 8 = 64$ token $\Rightarrow$ Self-Attention chi phí $64^2 = 4{,}096$ phép nhân -- rất nhỏ, nhưng mang lại global context cần thiết.
    
    \item Stage 2 có nhiều block nhất (6 blocks cho SA12), chiếm phần lớn tính toán -- tương tự ``stage 3'' của ResNet-50 chiếm phần lớn FLOPs.
\end{itemize}

% ----------------------------------------------------------------------------
\subsection{Downsampling Layers}
\label{subsec:downsampling}

Giữa các stage, \fastvit{} sử dụng PatchEmbed để giảm kích thước không gian $\times 2$ và tăng số kênh. PatchEmbed gồm:

\begin{enumerate}
    \item \textbf{ReparamLargeKernelConv ($7 \times 7$, depthwise, stride 2):} Tích chập nhân lớn giúp trích xuất đặc trưng phong phú tại ranh giới downsampling. Khi huấn luyện, gồm nhánh $7 \times 7$ và nhánh $3 \times 3$; khi suy luận, gộp thành một conv $7 \times 7$.
    
    \item \textbf{MobileOne Block ($1 \times 1$, groups=1):} Pointwise conv để ánh xạ kênh từ $C_i \to C_{i+1}$.
\end{enumerate}

Chi phí FLOPs của mỗi lớp downsampling (ví dụ $C_0=64 \to C_1=128$, resolution $64 \to 32$):
\begin{equation}
    \text{FLOPs}_{\text{down}} = 49 \times 64 \times 32^2 + 64 \times 128 \times 32^2 \approx 11.5\text{M FLOPs}
\end{equation}

% ----------------------------------------------------------------------------
\subsection{Classification Head}
\label{subsec:cls_head}

Classification Head gồm 3 bước:
\begin{enumerate}
    \item \textbf{MobileOne Block ($3 \times 3$, depthwise, SE):} Mở rộng kênh $C_3 \to C_3 \times r$ ($r = 2.0$ mặc định) với Squeeze-and-Excitation \cite{se2018} để tái chuẩn hóa đặc trưng kênh.
    \item \textbf{Global Average Pooling (GAP):} Pooling toàn cục $8 \times 8 \to 1 \times 1$.
    \item \textbf{Linear Layer:} Ánh xạ $C_3 \times r \to \text{num\_classes}$.
\end{enumerate}

% ============================================================================
\section{Structural Reparameterization trong FastViT}
\label{sec:reparameterization}

Structural Reparameterization là kỹ thuật then chốt giúp \fastvit{} đạt hiệu quả suy luận vượt trội. Nguyên tắc hoạt động:

\textbf{Khi huấn luyện (Training Mode):}
\begin{itemize}
    \item Sử dụng kiến trúc multi-branch với nhiều nhánh tích chập song song (conv $3 \times 3$ + BN, conv $1 \times 1$ + BN, identity + BN).
    \item Overparameterization giúp landscape loss mượt hơn, dễ tối ưu hơn, gradient flow tốt hơn.
\end{itemize}

\textbf{Khi suy luận (Inference Mode):}
\begin{itemize}
    \item Tất cả nhánh được hợp nhất thành \textbf{một lớp Conv duy nhất} thông qua:
    \begin{enumerate}
        \item \textbf{Conv-BN Fusion} (Phương trình~\ref{eq:conv_bn_fusion}): Hấp thụ BN vào trọng số Conv.
        \item \textbf{Branch Fusion} (Phương trình~\ref{eq:branch_fusion}): Cộng trọng số các nhánh (pad kernel nhỏ về cùng kích thước).
    \end{enumerate}
    \item Kết quả: mô hình suy luận chỉ gồm chuỗi Conv-Activation đơn giản, tối ưu cho phần cứng.
\end{itemize}

\textbf{Các thành phần sử dụng reparameterization trong \fastvit{}:}

\begin{table}[H]
\centering
\caption{Các module sử dụng Structural Reparameterization trong \fastvit{}.}
\label{tab:reparam_modules}
\begin{tabularx}{\textwidth}{l X X}
\toprule
\textbf{Module} & \textbf{Nhánh huấn luyện} & \textbf{Sau reparameterize} \\
\midrule
MobileOne Block & Conv $3 \times 3$ + BN, Conv $1 \times 1$ + BN, Identity + BN & Conv $3 \times 3$ đơn \\
RepMixer & Mixer (MobileOne) + Norm (MobileOne) + LayerScale & Conv $3 \times 3$ depthwise đơn \\
RepCPE & Conv depthwise + Identity & Conv depthwise đơn \\
ReparamLargeKernel & Conv $7 \times 7$ + BN, Conv $3 \times 3$ + BN & Conv $7 \times 7$ đơn \\
\bottomrule
\end{tabularx}
\end{table}

% ============================================================================
\section{RepMixer Block}
\label{sec:repmixer}

% ----------------------------------------------------------------------------
\subsection{Nguyên lý hoạt động}
\label{subsec:repmixer_principle}

RepMixer là thành phần token mixer cốt lõi của \fastvit{}, thay thế Self-Attention ở các stage đầu. RepMixer Block tuân theo cấu trúc MetaFormer \cite{metaformer2022}: Token Mixer + ConvFFN + Residual.

\textbf{Token Mixer (RepMixer):}

Khi huấn luyện:
\begin{equation}
    Y = X + \lambda \cdot \left(\text{Mixer}(X) - \text{Norm}(X)\right)
    \label{eq:repmixer_train}
\end{equation}

\noindent trong đó:
\begin{itemize}
    \item $X$: đầu vào, tensor $(B, C, H, W)$
    \item $\text{Mixer}$: MobileOne Block (depthwise conv $3 \times 3$, không activation)
    \item $\text{Norm}$: MobileOne Block (depthwise conv, chỉ có BN, không conv branch)
    \item $\lambda$: Layer Scale parameter (khởi tạo $10^{-5}$, có thể học)
\end{itemize}

Sau reparameterization, cả Mixer và Norm được rút gọn thành conv đơn, và toàn bộ biểu thức trở thành:
\begin{equation}
    W_{\text{final}} = I_{\text{conv}} + \lambda \cdot (W_{\text{mixer}} - W_{\text{norm}})
    \label{eq:repmixer_fused_w}
\end{equation}
\begin{equation}
    b_{\text{final}} = \lambda \cdot (b_{\text{mixer}} - b_{\text{norm}})
    \label{eq:repmixer_fused_b}
\end{equation}

\noindent Tức là: \textbf{một lớp depthwise conv $3 \times 3$ duy nhất} thực hiện toàn bộ token mixing + skip connection + layer scale.

\textbf{ConvFFN:}

Sau token mixer, ConvFFN mở rộng và nén đặc trưng kênh:
\begin{equation}
    \text{ConvFFN}(X) = \text{Conv}_{1\times1}^{(2)}\left(\text{GELU}\left(\text{Conv}_{1\times1}^{(1)}\left(\text{DWConv}_{7\times7}(X)\right)\right)\right)
\end{equation}

\noindent với DWConv $7 \times 7$ (depthwise) cung cấp thông tin cục bộ bổ sung, Conv $1 \times 1$ thứ nhất mở rộng $C \to C \times r$ ($r = 4$), Conv $1 \times 1$ thứ hai nén $C \times r \to C$.

\textbf{Tổng hợp RepMixer Block:}
\begin{equation}
    X' = \text{RepMixer}(X), \quad \text{Output} = X' + \text{DropPath}(\lambda' \cdot \text{ConvFFN}(X'))
\end{equation}

% ----------------------------------------------------------------------------
\subsection{So sánh với Self-Attention}
\label{subsec:repmixer_vs_attention}

\begin{table}[H]
\centering
\caption{So sánh chi phí tính toán RepMixer vs. Self-Attention (1 block, $C=256$, $k=3$).}
\label{tab:repmixer_vs_attention}
\begin{tabularx}{\textwidth}{X X X c}
\toprule
\textbf{Chỉ số} & \textbf{RepMixer (Inference)} & \textbf{Self-Attention} & \textbf{Tỷ lệ} \\
\midrule
Token mixer FLOPs ($N = 4096$) & $4096 \times 9 = 37$K & $4096^2 \times 256 = 4.3$G & $\sim 116{,}000\times$ \\
Token mixer FLOPs ($N = 1024$) & $1024 \times 9 = 9$K & $1024^2 \times 256 = 268$M & $\sim 29{,}000\times$ \\
Token mixer FLOPs ($N = 64$) & $64 \times 9 = 576$ & $64^2 \times 256 = 1$M & $\sim 1{,}800\times$ \\
Tham số (token mixer) & $9 \times C$ & $3 \times C^2 + C^2$ & $\sim C/2.25\times$ \\
Bộ nhớ trung gian & $\bigO{N \cdot C}$ & $\bigO{N^2 + N \cdot C}$ & -- \\
\bottomrule
\end{tabularx}
\end{table}

\textbf{Nhận xét quan trọng:} Với $N = 4{,}096$ token (stage 0), RepMixer nhanh hơn Self-Attention khoảng \textbf{100,000 lần} chỉ tính riêng phần token mixing. Đây là lý do \fastvit{} chỉ sử dụng attention ở stage cuối ($N = 64$) nơi chi phí attention nhỏ.

% ----------------------------------------------------------------------------
\subsection{Ưu điểm}
\label{subsec:repmixer_advantages}

\begin{enumerate}
    \item \textbf{Chi phí tuyến tính:} $\bigO{N \cdot k^2 \cdot C}$ thay vì $\bigO{N^2 \cdot C}$ -- scalable với ảnh phân giải cao.
    
    \item \textbf{Hiệu quả phần cứng:} Depthwise conv $3 \times 3$ được tối ưu cao trên mọi nền tảng (GPU, CPU, Mobile).
    
    \item \textbf{Reparameterizable:} Sau reparameterization, không có skip connection hay layer scale -- chỉ còn 1 phép conv, giảm memory access cost.
    
    \item \textbf{Overparameterized training:} Vẫn hưởng lợi từ multi-branch training để tăng accuracy.
    
    \item \textbf{Locality inductive bias:} Depthwise conv tự nhiên có tính cục bộ, phù hợp cho các đặc trưng low-level/mid-level.
\end{enumerate}

% ============================================================================
\section{Attention Block trong FastViT}
\label{sec:attention_block}

Ở stage cuối, \fastvit{} (các phiên bản SA/MA) sử dụng Attention Block với các thành phần:

\begin{enumerate}
    \item \textbf{RepCPE (Conditional Positional Encoding):} Nhúng thông tin vị trí không gian qua depthwise conv $7 \times 7$. Có thể reparameterize: $x + \text{PE}(x) \to \text{Conv}_{\text{fused}}(x)$ bằng cách cộng identity kernel vào trọng số PE.
    
    \item \textbf{MHSA (Multi-Headed Self-Attention):} Hỗ trợ đầu vào 4D ($B, C, H, W$). Flatten $\to$ attention $\to$ reshape lại. Head dim = 32 mặc định.
    
    \item \textbf{ConvFFN:} Tương tự RepMixer Block, sử dụng DWConv $7 \times 7$ + pointwise convs.
    
    \item \textbf{Layer Scale + DropPath:} Ổn định huấn luyện.
\end{enumerate}

Chi phí tính toán của Attention Block tại Stage 3 (SA12, $C = 512$, $N = 64$):
\begin{equation}
    \text{FLOPs}_{\text{attn}} = \underbrace{3 \times 64 \times 512^2}_{\text{QKV projection}} + \underbrace{64^2 \times 512}_{\text{Attention}} + \underbrace{64 \times 512^2}_{\text{Output proj}} \approx 53\text{M FLOPs/block}
\end{equation}

Với 2 blocks ở stage 3: $\sim 106$M FLOPs -- chỉ chiếm $\sim 5\%$ tổng FLOPs, nhưng mang lại global context quan trọng.

% ============================================================================
\section{Các phiên bản FastViT}
\label{sec:variants}

\fastvit{} cung cấp 7 phiên bản với quy mô và thiết kế khác nhau:

\begin{table}[H]
\centering
\caption{Cấu hình chi tiết các phiên bản \fastvit{}.}
\label{tab:variants}
\small
\begin{tabular}{l c c c c c c}
\toprule
\textbf{Phiên bản} & \textbf{Layers} & \textbf{Channels} & \textbf{MLP Ratio} & \textbf{Stage 3 Mixer} & \textbf{Params} & \textbf{Top-1} \\
\midrule
T8 & [2,2,4,2] & [48,96,192,384] & [3,3,3,3] & RepMixer & $\sim$4M & 76.2\% \\
T12 & [2,2,6,2] & [64,128,256,512] & [3,3,3,3] & RepMixer & $\sim$7M & 79.3\% \\
S12 & [2,2,6,2] & [64,128,256,512] & [4,4,4,4] & RepMixer & $\sim$9M & 79.9\% \\
SA12 & [2,2,6,2] & [64,128,256,512] & [4,4,4,4] & Attention & $\sim$11M & 80.9\% \\
SA24 & [4,4,12,4] & [64,128,256,512] & [4,4,4,4] & Attention & $\sim$21M & 82.7\% \\
SA36 & [6,6,18,6] & [64,128,256,512] & [4,4,4,4] & Attention & $\sim$31M & 83.6\% \\
MA36 & [6,6,18,6] & [76,152,304,608] & [4,4,4,4] & Attention & $\sim$44M & 83.9\% \\
\bottomrule
\end{tabular}
\end{table}

\textbf{Quy ước đặt tên:}
\begin{itemize}
    \item \textbf{T} (Tiny): Kênh nhỏ, MLP ratio 3. Nhẹ nhất, nhanh nhất.
    \item \textbf{S} (Small): Kênh tiêu chuẩn, MLP ratio 4. Cân bằng.
    \item \textbf{SA} (Small + Attention): Thêm Self-Attention ở stage cuối. Chính xác hơn S.
    \item \textbf{MA} (Medium + Attention): Kênh lớn hơn. Chính xác nhất.
    \item Số ($8, 12, 24, 36$): Tổng số blocks ($= \sum L_i$).
\end{itemize}

\textbf{Nhận xét về thiết kế:}
\begin{itemize}
    \item Từ S12 $\to$ SA12: thêm Attention ở stage 3, tăng 1\% accuracy chỉ với $\sim$2M params thêm.
    \item Từ SA12 $\to$ SA24: gấp đôi blocks, tăng 1.8\% accuracy.
    \item Từ SA36 $\to$ MA36: tăng channels (512 $\to$ 608), thêm 0.3\% accuracy nhưng thêm 13M params.
\end{itemize}

% ============================================================================
\section{Ưu điểm của FastViT}
\label{sec:fastvit_advantages}

\begin{enumerate}
    \item \textbf{Tốc độ suy luận vượt trội:} Sau reparameterization, mô hình chỉ gồm Conv + Activation, không có skip connection ở RepMixer, tối ưu cho mọi phần cứng.
    
    \item \textbf{Hiệu quả tham số:} SA12 đạt 80.9\% Top-1 với chỉ 11M params -- tương đương ViT-B/16 (86M params) kém hơn chỉ 4\%.
    
    \item \textbf{Scalable:} Dải rộng từ 4M (T8) đến 44M (MA36) params, phù hợp nhiều kịch bản triển khai.
    
    \item \textbf{Thiết kế lai thông minh:} CNN cho high-resolution (nhanh), Transformer cho low-resolution (global context).
    
    \item \textbf{Đa nhiệm:} Hỗ trợ cả classification, detection (qua \texttt{fork\_feat}) và segmentation.
    
    \item \textbf{Tương thích nền tảng:} Xuất CoreML, TorchScript, ONNX; hỗ trợ quantization INT8.
\end{enumerate}
